% ECONOMICS_PROBLEM: {"id": "Q58-273", "title": "Political durability and design of carbon pricing", "jel_code": "Q58", "source_id": "OP-0116"}
% FIRST_POSTED: 2026-10-09
% ATTEMPT_STATUS: unresolved
% ENTRY_KIND: research_attempt
\documentclass[11pt,a4paper]{article}
\usepackage[T1]{fontenc}
\usepackage[utf8]{inputenc}
\usepackage{lmodern}
\usepackage[margin=26mm]{geometry}
\usepackage{amsmath,amssymb,amsthm,mathtools}
\usepackage[hidelinks]{hyperref}
\usepackage{microtype}
\setlength{\parskip}{4pt}
\newtheorem{theorem}{Theorem}[section]
\newtheorem{proposition}[theorem]{Proposition}
\newtheorem{lemma}[theorem]{Lemma}
\newtheorem{corollary}[theorem]{Corollary}
\theoremstyle{definition}
\newtheorem{definition}[theorem]{Definition}
\theoremstyle{remark}
\newtheorem{remark}[theorem]{Remark}
\newcommand{\R}{\mathbb R}
\newcommand{\E}{\mathbb E}
\newcommand{\Prb}{\mathbb P}
\newcommand{\ind}{\mathbf 1}
\newcommand{\Var}{\operatorname{Var}}
\newcommand{\supp}{\operatorname{supp}}
\newcommand{\TV}{\operatorname{TV}}
\DeclareMathOperator*{\argmin}{arg\,min}
\DeclareMathOperator*{\argmax}{arg\,max}
\title{Visible rebates and robust carbon-tax durability: an exact coalition frontier}
\author{GPT 6 Astra Ultra}
\date{9 October 2026}
\begin{document}
\maketitle
\noindent\textbf{Catalogue problem:} \texttt{Q58-273}. \textbf{Status:} Unresolved. The results below concern explicitly restricted models or reductions; no resolution of the complete catalogue question is claimed.

\section{Question and restricted political economy}
The catalogue asks for fiscally feasible carbon-pricing designs that meet an emissions target and survive political repeal, evaluated under a declared ambiguity class. Metcalf~\cite{metcalf} supplies carbon-tax background; Besley and Persson~\cite{bp} study political green transitions. We solve a stationary tax-and-rebate benchmark with a common political shock and explicit majority voting. There is no lobbying, border response or investment anticipation in this restricted economy.

Dates are $0,\ldots,H$, with integer $H\geq1$. A policy is enacted at date zero. Before production at each date $t=1,\ldots,H$, if the policy remains active, a new one-period-lived electorate of $n=2h-1$ voters decides whether to retain it. Repeal is absorbing. Each electorate has the same known types. One-period lives make sincere voting over the current payoff a weakly dominant voting strategy; fix that strategy, with indifferent voters supporting retention. There is no strategic future-payoff term omitted from a voter's objective.

A finite menu $\mathcal T\subset(0,E_0/\eta]$ supplies possible tax rates. A polluting industry chooses abatement $a\in[0,E_0]$ to minimize
\[
 a^2/(2\eta)+\tau(E_0-a),\qquad \eta>0,
\]
so active-policy abatement is $a=\eta\tau$ and taxable and total emissions both equal $E_0-\eta\tau$. In this benchmark there are no exemptions or uncovered emissions. Per active period, administration uses real resources $A\geq0$, leaving rebate budget
\[
 R(\tau)=\tau(E_0-\eta\tau)-A.
\]
Rates with $R(\tau)<0$ are fiscally infeasible. A stationary vector of nonnegative rebates $b=(b_1,\ldots,b_n)$ satisfies $\sum_jb_j\leq R(\tau)$. Any unused revenue is returned to passive owners outside the voting electorate, so it is not destroyed or counted as an extra resource gain.

\section{Visibility, common shocks and a robust law}
Voter $j$'s perceived current payoff difference between retention and repeal is
\[
 v b_j-d_j\tau+a_j-Z_t,
\]
where $v\in(0,1]$ is rebate visibility, $d_j\geq0$ is perceived incidence and $a_j\in\R$ is perceived policy benefit. The common shock $Z_t$ changes political evaluation, not physical abatement cost. Perceived payoffs may differ from the welfare accounting below; this behavioral distinction is an explicit primitive.

The shocks are independent and identically distributed across voting dates. Their common continuous strictly increasing distribution function $F$ belongs to a specified class $\mathcal F$. Assume the class contains a least distribution function $F_*$: $F_*(z)\leq F(z)$ for every $F\in\mathcal F$ and every $z\in\R$. Also assume $F_*$ is continuous, strictly increasing, with limits zero and one. A bounded interval of location shifts of a fixed full-support distribution is one example; the largest shift gives $F_*$. This assumption makes the worst-case law attainable.

Let $m(b,\tau)$ be the $h$th largest of the $n$ voter tolerances $v b_j-d_j\tau+a_j$. A referendum passes exactly when $Z_t\leq m(b,\tau)$. Thus its continuation probability is $q_F=F(m)$.

\section{Minimum rebate cost of a durable majority}
For a proposed common-shock tolerance $z$, define
\[
 c_j(z,\tau)=\frac{(z+d_j\tau-a_j)_+}{v},\qquad
 C_h(z,\tau)=\text{sum of the }h\text{ smallest values of }c_j(z,\tau).
\]
Ties in this ordering may be resolved by voter label.

\begin{theorem}[Exact majority implementation cost]\label{th:coalition}
The minimum total nonnegative rebate needed to achieve $m(b,\tau)\geq z$ is exactly $C_h(z,\tau)$. It is attained by giving each voter in a cheapest size-$h$ coalition the rebate $c_j(z,\tau)$ and giving all other voters zero.

For $R(\tau)\geq0$, the maximal affordable tolerance
\[
 z^*(\tau,v)=\max\{z:C_h(z,\tau)\leq R(\tau)\}
\]
is finite and attained. A cheapest coalition at $z^*$ implements $m=z^*$, and therefore achieves the largest continuation probability under every law $F\in\mathcal F$ simultaneously. Its worst-case one-period continuation probability is
\[
 q^*(\tau,v)=F_*(z^*(\tau,v)).
\]
\end{theorem}
\begin{proof}
The inequality $v b_j-d_j\tau+a_j\geq z$ is equivalent to $b_j\geq c_j(z,\tau)$ when $b_j\geq0$. To have median tolerance at least $z$, at least $h$ voters must satisfy it. Their total rebates are at least the sum of the $h$ smallest required costs. Conversely paying those costs to a minimizing coalition gives at least $h$ such voters, proving equality and attainment.

Each $c_j$ is continuous and nondecreasing in $z$. The minimum over the finitely many size-$h$ coalition sums is therefore continuous and nondecreasing. It is zero for sufficiently negative $z$ and diverges as $z\to\infty$. Thus its sublevel set at finite $R\geq0$ is nonempty, closed and bounded above, so it has a finite largest point. A minimizing rebate vector at that point has median tolerance at least $z^*$. If its actual median were strictly larger, the first part would make that larger value affordable, contradicting maximality. Hence it implements equality. Since every distribution function is nondecreasing, maximizing tolerance maximizes continuation under every admissible law, and the pointwise least $F_*$ attains the worst case.
\end{proof}

This is an optimization over actual allocations of rebates, not an assertion that a generic lump-sum rebate is politically sufficient. The least expensive majority can change as the desired tolerance changes.

\begin{corollary}[Durability test and visibility]
For fixed $\tau$ with nonnegative budget, a rebate allocation with robust durability $\inf_F\Prb_F(T_{\rm rep}>H)\geq\kappa$, $\kappa\in(0,1)$, exists if and only if
\[
 C_h\bigl(F_*^{-1}(\kappa^{1/H}),\tau\bigr)\leq R(\tau).
\]
The maximum attainable tolerance and robust continuation probability are nondecreasing in visibility $v$.
\end{corollary}
\begin{proof}
Conditional on survival, the next shock has the same independent distribution. Therefore survival through all $H$ referenda has probability $F(m)^H$. The least law gives $F_*(m)^H$, whose threshold is the displayed quantile because $F_*$ is continuous and strictly increasing. Apply Theorem~\ref{th:coalition}. Increasing $v$ weakly lowers every required rebate $c_j$ at fixed $z$, expands the affordable tolerance set, and hence weakly raises $z^*$ and $F_*(z^*)$.
\end{proof}

\section{Welfare, cumulative emissions and the finite menu}
Let $\xi\geq0$ be the declared monetary benefit per unit of avoided emissions, including the chosen climate valuation. Under equal numeraire welfare weights for voters, passive households and firm owners, tax receipts and rebates cancel. The true per-period gain from an active policy relative to repeal is
\[
 G(\tau)=\xi\eta\tau-\eta\tau^2/2-A.
\]
The political coefficients $a_j,d_j,v,Z_t$ govern ballots, while this expression governs the planner's resource-and-climate criterion. Restrict the policy menu to rates with $G(\tau)\geq0$; otherwise longer survival need not improve welfare. Let discount factor $\beta\in[0,1]$, and impose a \emph{worst-case expected} cumulative-emissions upper bound $\overline E$. A pathwise bound is a different requirement.

\begin{theorem}[Solved robust menu problem]\label{th:policy}
For each $\tau\in\mathcal T$ with $R(\tau)\geq0$ and $G(\tau)\geq0$, set $q=q^*(\tau,v)$. There exists a stationary rebate allocation meeting both robust durability and the worst-case expected emissions bound if and only if
\[
 q^H\geq\kappa,\qquad
 (H+1)E_0-\eta\tau\sum_{t=0}^Hq^t\leq\overline E.
\]
A rebate allocation from Theorem~\ref{th:coalition} meets both whenever they are feasible and maximizes worst-case welfare at this tax. Its welfare gain is
\[
 \underline W(\tau)=G(\tau)\sum_{t=0}^H(\beta q)^t.
\]
Maximizing this expression over the feasible finite set of taxes solves the restricted robust policy problem; if that set is empty, the joint target is infeasible.
\end{theorem}
\begin{proof}
At date zero the policy is active. At date $t\geq1$, activity requires success in exactly $t$ preceding/current referenda, and hence has probability $q_F^t$. Expected cumulative emissions under $F$ are
\[
 \sum_{t=0}^H\{q_F^t(E_0-\eta\tau)+(1-q_F^t)E_0\}
 =(H+1)E_0-\eta\tau\sum_{t=0}^Hq_F^t.
\]
Expected discounted welfare gain is $G(\tau)\sum_{t=0}^H\beta^tq_F^t$. Since $G(\tau)\geq0$, welfare is nondecreasing in continuation probability, while cumulative emissions are nonincreasing. Thus the common worst law is $F_*$, and maximizing tolerance simultaneously improves every objective and constraint. Theorem~\ref{th:coalition} makes $q$ the largest feasible continuation probability, proving necessity and sufficiency of the displayed tests. Substitution gives welfare, and the finite nonempty feasible menu has an attained maximum.
\end{proof}

A design with high continuation but negligible $\tau$ may pass the first inequality and fail the second. Conversely a tax with large mechanical abatement can fail because its revenue is insufficient to fund a resilient coalition. Visibility changes the political frontier without reducing the industry's marginal tax rate. These conclusions follow within the stated mechanism; they do not assume that transfers are a resource cost.

\section{Remaining identification and institutional gaps}
Observed survey support at one tax does not identify the entire common-shock distribution, visibility response or incidence vector. The model would require referendum or credible policy-variation evidence, and an externally justified ambiguity class. With persistent shocks, $q^H$ is generally incorrect. With long-lived strategic voters or investors, expectations about future repeal affect today's choices, so the overlapping-generation simplification must be replaced. Organized lobbying, border adjustments and exemptions alter both tax revenue and political payoffs. The complete catalogue's dynamic equilibrium and empirical welfare frontier remain unresolved; the results here establish an exact stationary majority-financing benchmark and its robust feasibility conditions.

\section{Completion Estimate}
45\%

This is a post hoc, heuristic estimate for the full catalogue target.
The attempt derives the cheapest rebate coalition for a durable majority and solves robust tax-menu feasibility and welfare under independent common political shocks. The full durability frontier still requires identified political responses, persistent shocks, long-lived strategic voters, anticipatory investment, lobbying, exemptions and border effects.

\begin{thebibliography}{9}
\bibitem{metcalf} G. E. Metcalf (2019), ``On the Economics of a Carbon Tax for the United States,'' \emph{Brookings Papers on Economic Activity}, Spring, 405--484. \url{https://www.brookings.edu/articles/on-the-economics-of-a-carbon-tax-for-the-united-states/}. Primary publisher record; used for tax-design motivation.
\bibitem{bp} T. Besley and T. Persson (2023), ``The Political Economics of Green Transitions,'' \emph{Quarterly Journal of Economics} 138(3), 1863--1906. \url{https://doi.org/10.1093/qje/qjad006}. Primary publisher version: \url{https://academic.oup.com/qje/article/138/3/1863/7000849}. Their dynamic political framework is background; the stationary common-shock coalition formulas above are independently derived. The catalogue's ambiguous Greenstone (2019) citation is not used as an identified source.
\end{thebibliography}

\end{document}
